\chapter{Literature Review}
\label{chap:literature_review}

This chapter sets the physical and computational context for the project. It
reviews the Standard Model and the role of the top quark, the formation and
substructure of jets, the engineered observables used in classical tagging,
and the deep-learning architectures that define the current state of the art.

\section{The Standard Model and the Top Quark}
\label{sec:sm}

The Standard Model (SM) of particle physics describes the electromagnetic,
weak and strong interactions of quarks and leptons through gauge symmetry
$SU(3)_C\times SU(2)_L\times U(1)_Y$~\citep{pdg2022}. The top quark, the
up-type quark of the third generation, is the heaviest known elementary
particle, $m_t\simeq172.76$\,GeV. Because of this large mass its lifetime is
shorter than the timescale of hadronisation, so it decays as a bare quark,
almost always via $t\rightarrow Wb$. The $W$ boson subsequently decays either
leptonically ($W\rightarrow\ell\nu$) or hadronically
($W\rightarrow q\bar{q}'$). The hadronic channel,
$t\rightarrow Wb\rightarrow q\bar{q}'b$, is the one relevant to this thesis:
it produces three quarks that, for a boosted top, merge into one
large-radius jet.

\section{Jets and Jet Reconstruction}
\label{sec:jets}

Coloured partons cannot exist freely; they fragment and hadronise into
collimated streams of colourless hadrons. These streams are reconstructed
into jets by sequential-recombination algorithms, most commonly the
anti-$k_t$ algorithm~\citep{cacciari2008} as implemented in
\textsc{FastJet}~\citep{cacciari2012fastjet}. For boosted-object studies a
large radius parameter ($R=0.8$ or $1.0$) is used so that the entire decay is
contained in a single jet. A boosted top jet then differs from an ordinary
QCD jet in three observable ways: it is more massive (its mass peaks near
$m_t$), it has a three-pronged rather than one-pronged energy flow, and it
contains more constituents. Exploiting these differences is the essence of
jet-substructure analysis~\citep{larkoski2020,kogler2019}.

\section{Engineered Substructure Observables}
\label{sec:observables}

Classical taggers compress the jet into a handful of physically motivated
numbers. The most important are:
\begin{itemize}
    \item \textbf{Jet mass} --- the invariant mass of the summed constituent
    four-momenta, which discriminates the heavy top from light QCD jets.
    \item \textbf{$N$-subjettiness} $\tau_N$ and its ratios
    $\tau_{21}=\tau_2/\tau_1$, $\tau_{32}=\tau_3/\tau_2$, which quantify how
    well the radiation is described by $N$ subjet axes~\citep{thaler2011}. A
    genuine three-prong top jet has small $\tau_{32}$.
    \item \textbf{Energy-correlation functions} and the ratios $C_2$, $D_2$,
    which capture $n$-point energy correlations without explicit subjet
    finding~\citep{larkoski2013,larkoski2014}.
\end{itemize}
Fed to a boosted decision tree (BDT), typically implemented with
\textsc{XGBoost}~\citep{chen2016xgboost}, these observables already provide
strong discrimination and remain a competitive, interpretable baseline.

\section{Deep Learning on Jets}
\label{sec:deep}

The deep-learning era of tagging replaced hand-engineered features with
representations learned directly from the constituents. Three families are
relevant here.

\subsection{Jet Images and Convolutional Networks}
Pixelating the constituents in the rapidity--azimuth plane turns a jet into
an image that a convolutional neural network (CNN) can classify, in close
analogy with calorimeter readout~\citep{deoliveira2016,komiske2017}. Image
methods were the first deep approaches to surpass engineered-feature BDTs,
but the binning discards angular resolution and breaks the natural
permutation symmetry of the constituents.

\subsection{Particle Clouds and Graph Networks}
Treating a jet as an unordered set --- a ``particle cloud'' --- preserves the
full constituent information. ParticleNet~\citep{qu2020particlenet} adapts the
Dynamic Graph CNN~\citep{wang2019dgcnn} to jets, building a
$k$-nearest-neighbour graph in feature space and applying EdgeConv operations
that are invariant to constituent ordering. This and related set/graph
approaches~\citep{komiske2019efn,moreno2020jedinet} substantially improved
tagging performance.

\subsection{Attention and the Particle Transformer}
The Transformer architecture~\citep{vaswani2017attention} models pairwise
interactions through self-attention. The Particle
Transformer~\citep{qu2022particletransformer} augments standard attention
with physically motivated pairwise features (angular distance, $k_t$,
momentum fraction and invariant mass) injected as an attention bias, and
currently holds the top of the reference-dataset
leaderboard~\citep{kasieczka2019landscape}. Related physics-informed and
Lorentz-equivariant designs continue to push the
frontier~\citep{bogatskiy2020lorentz,gong2022lorentznet}.

\section{Probability Calibration in Machine Learning}
\label{sec:calibration_lit}

A classifier is \emph{calibrated} if, among samples it assigns score $p$, a
fraction $p$ are truly positive. \citet{guo2017calibration} showed that
modern deep networks, although accurate, are frequently \emph{over-confident}
and that a simple post-hoc \emph{temperature scaling} --- dividing the logits
by a single learned scalar $T$ --- often restores calibration. Calibration is
measured by the Expected Calibration Error (ECE) and visualised with
reliability diagrams~\citep{niculescu2005calibration}. Although calibration
is standard in the broader ML community, it is seldom reported in
physics-tagging papers, which is precisely the gap this thesis explores.

\section{Summary and Gap}
\label{sec:gap}

The literature establishes both a strong set of engineered observables and a
clear hierarchy of deep models, all benchmarked at the full data and compute
budget. What is missing, and what this thesis provides, is a controlled,
reproducible comparison of \emph{data efficiency} and \emph{probability
calibration} across these architectures in the modest-budget regime that is
realistic for most students and small groups.
